A study from the Human Brain Project provides a compelling mechanism for \textit{how} the brain might generate this functional, bizarre data, drawing inspiration from machine learning \citep{franceschi2024gans}. The model uses a concept from AI called Generative Adversarial Networks (GANs), in which two neural networks compete to generate novel, realistic data from an existing dataset.

This study proposes a complementary, multi-stage learning algorithm for sleep:
\begin{enumerate}
    \item \textbf{Wakefulness:} The brain is exposed to new stimuli and acquires data (e.g., sees pictures of cars).
    \item \textbf{NREM Sleep:} The brain "replays" and "solidifies" this experience, stabilizing the memory.
    \item \textbf{REM Sleep:} The brain's "GAN" mechanism activates, generating "twisted but realistic versions and combinations" of the cars. This is the "bizarre" dream, which functions as adversarial training data.
\end{enumerate}

The functional proof of this model is
that when the researchers \textit{suppressed} this "bizarre" REM phase, the network's learning accuracy and, most importantly, its ability to "discover semantic concepts" (to generalize) \textit{decreased} \citep{franceschi2024gans}. The conclusion is that "this bizarreness serves a purpose," allowing the brain to move from rote memorization to true conceptual understanding.